% =====================================================================
% Rozdziały pracy magisterskiej -- prezentacja wyników eksperymentów
%
% Plik zawiera trzy sekcje przedstawiające wyniki badań nad detekcją
% fałszywych informacji (fake news) z wykorzystaniem trzech rodzin
% modeli: klasycznych modeli uczenia maszynowego, sieci głębokich oraz
% modeli opartych na architekturze Transformer.
%
% Wyniki są oparte na czterech zbiorach wyników eksperymentów:
%   * results_test.csv          -- ewaluacja na zbiorze testowym
%                                  (split 10\% z każdego korpusu),
%   * results_text_web.csv      -- ewaluacja na pełnych tekstach
%                                  artykułów pozyskanych ze stron WWW,
%   * results_title_web.csv     -- ewaluacja na samych tytułach
%                                  artykułów ze stron WWW,
%   * results_short_text_web.csv-- ewaluacja na krótkich fragmentach
%                                  tekstu (lead/podsumowanie).
%
% Aby skompilować plik wymagane są pakiety:
%   \usepackage[utf8]{inputenc}
%   \usepackage[T1]{fontenc}
%   \usepackage[polish]{babel}
%   \usepackage{booktabs}
%   \usepackage{array}
%   \usepackage{multirow}
%   \usepackage{caption}
%   \usepackage{graphicx}
% =====================================================================

\section{Wyniki klasycznych modeli uczenia maszynowego}
\label{sec:wyniki_klasyczne}

W~niniejszym rozdziale przedstawione zostały wyniki eksperymentów
przeprowadzonych przy użyciu ośmiu klasycznych algorytmów uczenia
maszynowego: regresji logistycznej (\textit{lr}), maszyny wektorów
nośnych (\textit{svm}), drzewa decyzyjnego (\textit{dt}), lasu
losowego (\textit{rf}), gradient boostingu (\textit{xgb}), naiwnego
klasyfikatora Bayesa w~wariantach Gaussa (\textit{nb}) i~wielomianowym
(\textit{mnb}) oraz $k$-najbliższych sąsiadów (\textit{knn}). Każdy
klasyfikator został przetestowany w~kombinacji z~czterema reprezentacjami
tekstu: TF--IDF, bag-of-words (BoW), Word2Vec oraz GloVe. Niedozwolone
kombinacje (np.\ wielomianowy Bayes z~uśrednionymi gęstymi wektorami)
były automatycznie odrzucane przez fabrykę modeli, dlatego liczność
eksperymentów dla każdego korpusu nie jest identyczna w~każdym wierszu.
Hiperparametry zostały dobrane za pomocą biblioteki Optuna
(samplerem był \textit{TPESampler}), a~ostateczne metryki uzyskano na
niezależnym podzbiorze testowym (10\,\% każdego ze zbiorów: ISOT, LIAR
oraz WELFake).

\subsection{Wyniki na zbiorach testowych}
\label{sec:wyniki_klasyczne_test}

Tabela~\ref{tab:klas_isot} zawiera podsumowanie najważniejszych
wyników klasycznych modeli na zbiorze ISOT. Wartości w~tabeli zostały
posortowane malejąco względem miary $F_1$.

\begin{table}[ht]
\centering
\caption{Wyniki klasycznych modeli uczenia maszynowego na zbiorze
testowym ISOT (najlepsze osiem kombinacji wg miary $F_1$).}
\label{tab:klas_isot}
\begin{tabular}{llrrrrrr}
\toprule
Model & Embedding & Acc. & Prec. & Recall & $F_1$ & $t_{tr}$ [s] & $t_{inf}$ [s] \\
\midrule
xgb & tfidf    & 0,9977 & 0,9977 & 0,9977 & 0,9977 & 150,05 & 0,072 \\
xgb & bow      & 0,9974 & 0,9975 & 0,9974 & 0,9974 & 12,14  & 0,063 \\
lr  & tfidf    & 0,9921 & 0,9921 & 0,9921 & 0,9921 & 1,77   & 0,040 \\
svm & tfidf    & 0,9905 & 0,9905 & 0,9905 & 0,9905 & 29,39  & 0,048 \\
svm & word2vec & 0,9836 & 0,9836 & 0,9836 & 0,9836 & 6,60   & 0,023 \\
lr  & word2vec & 0,9813 & 0,9813 & 0,9813 & 0,9813 & 3,30   & 0,016 \\
rf  & tfidf    & 0,9739 & 0,9743 & 0,9739 & 0,9739 & 27,67  & 0,607 \\
mnb & tfidf    & 0,9619 & 0,9619 & 0,9619 & 0,9619 & 0,12   & 0,034 \\
\bottomrule
\end{tabular}
\end{table}

Na zbiorze ISOT klasyczne modele osiągnęły bardzo wysoką jakość
klasyfikacji. Najlepszy wynik (\(F_1 = 0{,}9977\)) uzyskał XGBoost
z~reprezentacją TF--IDF, jednak różnica między pierwszą a~ósmą pozycją
rankingu wynosi niespełna 4 punkty procentowe. Warto zauważyć, że
reprezentacje rzadkie (TF--IDF, BoW) wyraźnie dominują nad gęstymi
uśrednionymi wektorami (Word2Vec, GloVe) -- pierwsze cztery miejsca
zajmują wyłącznie modele oparte na TF--IDF lub BoW. Najsłabszy z~modeli,
$k$-NN z~reprezentacją BoW, uzyskał jedynie \(F_1 = 0{,}6444\), co
wynika z~problemu rzadkości i~wysokiej wymiarowości przestrzeni cech,
w~której metryka odległości euklidesowej traci interpretowalność.

Tabela~\ref{tab:klas_welfake} zawiera analogiczne wyniki dla korpusu
WELFake.

\begin{table}[ht]
\centering
\caption{Wyniki klasycznych modeli na zbiorze testowym WELFake
(najlepsze osiem kombinacji wg miary $F_1$).}
\label{tab:klas_welfake}
\begin{tabular}{llrrrrrr}
\toprule
Model & Embedding & Acc. & Prec. & Recall & $F_1$ & $t_{tr}$ [s] & $t_{inf}$ [s] \\
\midrule
xgb & bow      & 0,9713 & 0,9715 & 0,9713 & 0,9713 & 56,17  & 0,173 \\
xgb & tfidf    & 0,9711 & 0,9713 & 0,9711 & 0,9711 & 147,67 & 0,045 \\
lr  & tfidf    & 0,9576 & 0,9576 & 0,9576 & 0,9576 & 8,14   & 0,072 \\
svm & tfidf    & 0,9490 & 0,9490 & 0,9490 & 0,9489 & 58,60  & 0,085 \\
dt  & bow      & 0,9362 & 0,9381 & 0,9362 & 0,9364 & 20,21  & 0,021 \\
dt  & tfidf    & 0,9336 & 0,9349 & 0,9336 & 0,9337 & 57,23  & 0,029 \\
svm & word2vec & 0,9194 & 0,9196 & 0,9194 & 0,9195 & 32,06  & 0,065 \\
lr  & word2vec & 0,9165 & 0,9165 & 0,9165 & 0,9165 & 5,55   & 0,016 \\
\bottomrule
\end{tabular}
\end{table}

Wyniki na korpusie WELFake są wyraźnie skromniejsze niż na ISOT, ale
nadal pozwalają na uznanie problemu za rozwiązany w~stopniu bardzo
dobrym. Najlepszą kombinacją okazał się XGBoost z~reprezentacją BoW
(\(F_1 = 0{,}9713\)), co potwierdza dużą skuteczność metod opartych
na drzewach decyzyjnych w~połączeniu z~rzadkimi reprezentacjami
tekstu. Średnia miara $F_1$ dla wszystkich klasycznych modeli na tym
zbiorze wyniosła \(0{,}883\), co świadczy o~większej trudności
zadania -- WELFake jest zbiorem zagregowanym z~czterech źródeł
i~zawiera bardziej zróżnicowane stylistycznie teksty.

Najtrudniejszy okazał się zbiór LIAR, który zawiera krótkie wypowiedzi
polityków sklasyfikowane na skali sześciostopniowej zbinaryzowanej do
problemu dwuklasowego (Tabela~\ref{tab:klas_liar}).

\begin{table}[ht]
\centering
\caption{Wyniki klasycznych modeli na zbiorze testowym LIAR
(najlepsze osiem kombinacji wg miary $F_1$).}
\label{tab:klas_liar}
\begin{tabular}{llrrrrrr}
\toprule
Model & Embedding & Acc. & Prec. & Recall & $F_1$ & $t_{tr}$ [s] & $t_{inf}$ [s] \\
\midrule
mnb & bow      & 0,6335 & 0,6321 & 0,6335 & 0,6326 & 0,01  & 0,008 \\
xgb & bow      & 0,6280 & 0,6315 & 0,6280 & 0,6291 & 3,81  & 0,025 \\
mnb & tfidf    & 0,6312 & 0,6274 & 0,6312 & 0,6232 & 0,01  & 0,013 \\
lr  & glove    & 0,6202 & 0,6280 & 0,6202 & 0,6215 & 1,47  & 0,015 \\
svm & glove    & 0,6186 & 0,6262 & 0,6186 & 0,6199 & 3,32  & 0,020 \\
rf  & tfidf    & 0,6108 & 0,6243 & 0,6108 & 0,6115 & 1,11  & 0,067 \\
nb  & glove    & 0,6092 & 0,6144 & 0,6092 & 0,6105 & 0,03  & 0,010 \\
xgb & tfidf    & 0,6085 & 0,6105 & 0,6085 & 0,6092 & 11,84 & 0,026 \\
\bottomrule
\end{tabular}
\end{table}

Na zbiorze LIAR żadna z~klasycznych metod nie przekroczyła progu
\(F_1 = 0{,}64\). Najlepszy wynik (\(F_1 = 0{,}6326\)) uzyskał
wielomianowy klasyfikator Bayesa z~reprezentacją BoW, co jest typowym
zachowaniem dla bardzo krótkich tekstów i~ograniczonego słownika.
Wynik bliski losowemu (baseline większościowy wynosi około
\(0{,}56\)) potwierdza, że proste cechy n-gramowe nie wystarczają,
aby uchwycić niuanse językowe charakteryzujące fałszywe wypowiedzi
polityków.

\subsection{Analiza czasów uczenia i~wnioskowania}
\label{sec:wyniki_klasyczne_czas}

Średnie czasy obliczeń dla rodziny klasycznej, zestawione w~podziale
na zbiory, podano w~Tabeli~\ref{tab:klas_czasy}.

\begin{table}[ht]
\centering
\caption{Średnie czasy uczenia i~wnioskowania klasycznych modeli
w~zależności od zbioru danych (w~sekundach).}
\label{tab:klas_czasy}
\begin{tabular}{lrrrr}
\toprule
Zbiór & $\overline{t_{tr}}$ & $\max{t_{tr}}$ & $\overline{t_{inf}}$ & $\max{t_{inf}}$ \\
\midrule
ISOT    & 18,02  & 150,05 & 0,755 & 11,02  \\
LIAR    & 3,21   & 16,38  & 0,044 & 0,234  \\
WELFake & 33,29  & 147,67 & 2,504 & 37,16  \\
\bottomrule
\end{tabular}
\end{table}

Czas uczenia dla klasycznych modeli pozostaje w~granicach kilkudziesięciu
sekund nawet dla największego zbioru. Najszybsze są klasyfikatory
naiwnego Bayesa oraz drzewa decyzyjne (rzędu setnych części sekundy),
najwolniejszy zaś XGBoost. Wnioskowanie jest zazwyczaj poniżej 0,1~s
na pełnym zbiorze testowym, z~wyjątkiem $k$-NN, którego złożoność rośnie
liniowo z~liczbą próbek (do \(37{,}16\)~s na WELFake). Daje to klasycznym
modelom przewagę pod kątem kosztu eksploatacyjnego w~porównaniu
z~modelami głębokimi i~transformerami (por.\ rozdz.~\ref{sec:wyniki_transformery}).

\subsection{Wyniki na danych pozyskanych ze stron WWW}
\label{sec:wyniki_klasyczne_web}

Modele wytrenowane na zbiorach ISOT, LIAR i~WELFake zostały następnie
poddane ewaluacji na zewnętrznym zbiorze \(1{,}5\)~tys.\ artykułów
pozyskanych metodą web-scrapingu z~popularnych portali informacyjnych.
Ewaluacja została przeprowadzona w~trzech wariantach: dla pełnego
tekstu, dla samych nagłówków oraz dla krótkich fragmentów (lead).
Tabela~\ref{tab:klas_web} zawiera najlepsze wyniki klasycznych modeli
w~każdym wariancie.

\begin{table}[ht]
\centering
\caption{Najlepsze klasyczne modele w~ewaluacji na danych WWW
(według miary $F_1$).}
\label{tab:klas_web}
\begin{tabular}{lllrrrr}
\toprule
Wariant & Zbiór & Model + Emb. & Acc. & Prec. & Recall & $F_1$ \\
\midrule
text       & LIAR    & dt + tfidf      & 0,8074 & 0,8210 & 0,8074 & 0,8140 \\
text       & ISOT    & nb + word2vec   & 0,6929 & 0,8969 & 0,6929 & 0,7533 \\
text       & WELFake & dt + bow        & 0,7136 & 0,7867 & 0,7136 & 0,7483 \\
title      & ISOT    & rf + bow        & 0,8799 & 0,8350 & 0,8799 & 0,8524 \\
title      & WELFake & rf + tfidf      & 0,8981 & 0,8066 & 0,8981 & 0,8499 \\
title      & LIAR    & svm + tfidf     & 0,6201 & 0,7800 & 0,6201 & 0,6934 \\
short text & ISOT    & rf + bow        & 0,8970 & 0,8684 & 0,8970 & 0,8747 \\
short text & WELFake & dt + bow        & 0,8839 & 0,8119 & 0,8839 & 0,8451 \\
short text & LIAR    & svm + tfidf     & 0,6554 & 0,7990 & 0,6554 & 0,7210 \\
\bottomrule
\end{tabular}
\end{table}

Z~przeprowadzonej analizy wynika, że klasyczne modele zaskakująco dobrze
generalizują na dane pozyskane spoza zbioru treningowego, szczególnie
gdy analizowany jest sam nagłówek lub krótki opis artykułu. Najlepsze
modele (\textit{rf}~+~\textit{bow} oraz \textit{dt}~+~\textit{bow}) na
nagłówkach i~lead-ach uzyskują \(F_1 \in [0{,}84;\,0{,}87]\), co
stanowi nieznaczny spadek względem zbiorów testowych in-domain.
Wyraźnie gorzej wypadają w~analizie pełnego tekstu artykułu, gdzie
średnia miara $F_1$ spada do okolic \(0{,}58\). Świadczy to o~tym,
że ich zdolność uogólniania ogranicza dystrybucyjne podobieństwo
między zbiorem treningowym a~testowym -- klasyczne reprezentacje
silnie zależą od częstotliwości występowania konkretnych słów.

\subsection{Podsumowanie}
\label{sec:wyniki_klasyczne_podsum}

Klasyczne modele uczenia maszynowego stanowią silny i~szybki punkt
odniesienia w~zadaniu detekcji fałszywych informacji. Na korpusie ISOT
najlepsze konfiguracje (XGBoost~+~TF-IDF) osiągają niemal idealną
skuteczność (\(F_1 = 0{,}9977\)), a~na WELFake skutecznie przekraczają
próg \(0{,}97\). Jednocześnie ich wydajność jest najwyższa w~całym
badaniu pod względem czasu uczenia i~wnioskowania, co czyni je
naturalnymi kandydatami w~systemach wdrożeniowych o~niskich opóźnieniach
(np.\ rozszerzenie do przeglądarki). Słabością tej rodziny jest zbiór
LIAR -- najtrudniejszy w~całym badaniu -- na którym żaden klasyczny
model nie przekroczył \(F_1 = 0{,}64\).

\section{Wyniki sieci głębokich}
\label{sec:wyniki_glebokie}

W~tym rozdziale omówione zostały wyniki czterech architektur sieci
głębokich: jednokierunkowej sieci LSTM, sieci GRU, dwukierunkowej
sieci LSTM (BiLSTM) oraz jednowymiarowej sieci konwolucyjnej (CNN).
Modele zostały zaimplementowane w~bibliotece PyTorch, a~ich
hiperparametry (rozmiar warstwy ukrytej, dropout, learning rate,
rozmiar partii) były dobierane przez Optuna. W~odróżnieniu od
klasycznych modeli sieci głębokie korzystają z~osadzeń słownych
wczytywanych z~zewnętrznych plików: Word2Vec (Google News, 300d)
oraz GloVe (6B, 100d).

\subsection{Wyniki na zbiorach testowych}
\label{sec:wyniki_glebokie_test}

Tabela~\ref{tab:dl_isot} zawiera podsumowanie wyników sieci głębokich
na zbiorze ISOT.

\begin{table}[ht]
\centering
\caption{Wyniki sieci głębokich na zbiorze testowym ISOT.}
\label{tab:dl_isot}
\begin{tabular}{llrrrrrr}
\toprule
Model & Embedding & Acc. & Prec. & Recall & $F_1$ & $t_{tr}$ [s] & $t_{inf}$ [s] \\
\midrule
gru    & word2vec & 0,9946 & 0,9946 & 0,9946 & 0,9946 & 67,06 & 0,491 \\
bilstm & word2vec & 0,9941 & 0,9941 & 0,9941 & 0,9941 & 71,13 & 0,388 \\
cnn    & word2vec & 0,9916 & 0,9916 & 0,9916 & 0,9916 & 83,04 & 0,425 \\
gru    & glove    & 0,9913 & 0,9913 & 0,9913 & 0,9913 & 55,28 & 0,675 \\
cnn    & glove    & 0,9880 & 0,9882 & 0,9880 & 0,9880 & 17,76 & 0,365 \\
bilstm & glove    & 0,9877 & 0,9879 & 0,9877 & 0,9877 & 39,21 & 0,633 \\
lstm   & glove    & 0,9854 & 0,9855 & 0,9854 & 0,9854 & 55,69 & 0,680 \\
lstm   & word2vec & 0,9844 & 0,9846 & 0,9844 & 0,9844 & 64,12 & 0,415 \\
\bottomrule
\end{tabular}
\end{table}

Wszystkie modele głębokie na zbiorze ISOT przekraczają granicę
\(F_1 = 0{,}98\), a~najlepsza konfiguracja (GRU~+~Word2Vec) osiąga
\(F_1 = 0{,}9946\). Średnia wartość miary $F_1$ wynosi \(0{,}9896\),
co jest najwyższym wynikiem wśród wszystkich rodzin modeli z~wyjątkiem
transformerów (por.\ Tabela~\ref{tab:bert_isot}). Można zauważyć,
że osadzenia Word2Vec systematycznie przewyższają GloVe -- pierwsze
trzy miejsca rankingu zajmują kombinacje z~tym pierwszym osadzeniem.
Architektury rekurencyjne (GRU, BiLSTM) okazują się minimalnie lepsze
od konwolucyjnych, co tłumaczyć można dłuższym kontekstem zdań
w~artykułach prasowych.

\begin{table}[ht]
\centering
\caption{Wyniki sieci głębokich na zbiorze testowym WELFake.}
\label{tab:dl_welfake}
\begin{tabular}{llrrrrrr}
\toprule
Model & Embedding & Acc. & Prec. & Recall & $F_1$ & $t_{tr}$ [s] & $t_{inf}$ [s] \\
\midrule
gru    & word2vec & 0,9824 & 0,9824 & 0,9824 & 0,9824 & 144,14 & 0,999 \\
bilstm & word2vec & 0,9812 & 0,9812 & 0,9812 & 0,9812 & 166,63 & 1,267 \\
cnn    & word2vec & 0,9788 & 0,9788 & 0,9788 & 0,9788 & 129,05 & 0,825 \\
gru    & glove    & 0,9783 & 0,9783 & 0,9783 & 0,9783 & 48,71  & 1,008 \\
cnn    & glove    & 0,9769 & 0,9771 & 0,9769 & 0,9769 & 36,25  & 0,798 \\
bilstm & glove    & 0,9761 & 0,9762 & 0,9761 & 0,9761 & 71,27  & 1,257 \\
lstm   & glove    & 0,9669 & 0,9670 & 0,9669 & 0,9669 & 49,79  & 1,023 \\
lstm   & word2vec & 0,9604 & 0,9605 & 0,9604 & 0,9604 & 193,55 & 1,306 \\
\bottomrule
\end{tabular}
\end{table}

Na korpusie WELFake sieci głębokie zachowują przewagę nad klasycznymi
modelami -- średnia miara $F_1$ wynosi \(0{,}9751\), co przekłada się
na zysk około 9~punktów procentowych względem rodziny klasycznej.
Najsłabszą architekturą pozostaje jednokierunkowy LSTM, co jest
spodziewane: brak informacji ,,z~przyszłości'' istotnie ogranicza
modelowanie struktur narracyjnych. Z~kolei BiLSTM oraz GRU osiągają
porównywalne wyniki przy wyraźnie krótszym czasie uczenia w~przypadku
GRU.

Diametralnie inaczej wyglądają wyniki na zbiorze LIAR
(Tabela~\ref{tab:dl_liar}).

\begin{table}[ht]
\centering
\caption{Wyniki sieci głębokich na zbiorze testowym LIAR.}
\label{tab:dl_liar}
\begin{tabular}{llrrrrr}
\toprule
Model & Embedding & Acc. & Prec. & Recall & $F_1$ & Uwagi \\
\midrule
bilstm & glove    & 0,6179 & 0,6152 & 0,6179 & 0,6157 & --- \\
cnn    & glove    & 0,6092 & 0,6040 & 0,6092 & 0,6006 & --- \\
cnn    & word2vec & 0,5638 & 0,5576 & 0,5638 & 0,5580 & --- \\
bilstm & word2vec & 0,5583 & 0,5795 & 0,5583 & 0,5562 & --- \\
gru    & word2vec & 0,5583 & 0,3117 & 0,5583 & 0,4001 & klasa większościowa \\
gru    & glove    & 0,4417 & 0,1951 & 0,4417 & 0,2706 & klasa większościowa \\
lstm   & word2vec & 0,4417 & 0,1951 & 0,4417 & 0,2706 & klasa większościowa \\
lstm   & glove    & 0,4417 & 0,1951 & 0,4417 & 0,2706 & klasa większościowa \\
\bottomrule
\end{tabular}
\end{table}

Wyniki sieci LSTM i~GRU na LIAR jednoznacznie wskazują na zjawisko
\textit{collapse} -- macierze pomyłek $\begin{pmatrix}564 & 0\\713 &
0\end{pmatrix}$ oraz $\begin{pmatrix}0 & 564\\0 & 713\end{pmatrix}$
oznaczają, że sieci przewidziały wyłącznie jedną etykietę dla
wszystkich próbek. Wynika to z~dwóch przyczyn: bardzo krótkich
sekwencji (mediana długości około 18 słów) oraz niewielkiego rozmiaru
zbioru treningowego (10\,269 wypowiedzi), co dla głębokich sieci
rekurencyjnych jest dalece niewystarczające. Jedynie architektury
BiLSTM (z~GloVe) oraz CNN potrafią pokonać granicę większościową --
najlepszy z~modeli, BiLSTM~+~GloVe, uzyskał \(F_1 = 0{,}6157\), co
jest wynikiem zbliżonym do najlepszych klasycznych modeli na tym
zbiorze.

\subsection{Analiza czasów uczenia i~wnioskowania}
\label{sec:wyniki_glebokie_czas}

\begin{table}[ht]
\centering
\caption{Średnie czasy uczenia i~wnioskowania sieci głębokich
w~zależności od zbioru danych (w~sekundach).}
\label{tab:dl_czasy}
\begin{tabular}{lrrrr}
\toprule
Zbiór & $\overline{t_{tr}}$ & $\max{t_{tr}}$ & $\overline{t_{inf}}$ & $\max{t_{inf}}$ \\
\midrule
ISOT    & 56,66  & 83,04  & 0,509 & 0,680 \\
LIAR    & 9,43   & 29,45  & 0,079 & 0,220 \\
WELFake & 104,92 & 193,55 & 1,060 & 1,306 \\
\bottomrule
\end{tabular}
\end{table}

Czasy uczenia sieci głębokich są przeciętnie 3--5 razy dłuższe od
klasycznych modeli (np.\ średnio \(104{,}9\)~s vs.\ \(33{,}3\)~s dla
WELFake), ale wciąż pozostają o~rząd wielkości krótsze od czasów
transformerów. Inferencja na pełnym zbiorze testowym mieści się
w~granicach jednej sekundy, co czyni je akceptowalnym wyborem dla
zastosowań czasu rzeczywistego.

\subsection{Wyniki na danych pozyskanych ze stron WWW}
\label{sec:wyniki_glebokie_web}

\begin{table}[ht]
\centering
\caption{Najlepsze sieci głębokie w~ewaluacji na danych WWW
(według miary $F_1$).}
\label{tab:dl_web}
\begin{tabular}{lllrrrr}
\toprule
Wariant & Zbiór & Model + Emb. & Acc. & Prec. & Recall & $F_1$ \\
\midrule
text       & LIAR    & bilstm + glove    & 0,8843 & 0,8388 & 0,8843 & 0,8550 \\
text       & WELFake & bilstm + word2vec & 0,5935 & 0,8073 & 0,5935 & 0,6732 \\
text       & ISOT    & cnn + word2vec    & 0,5879 & 0,8826 & 0,5879 & 0,6671 \\
title      & LIAR    & gru + word2vec    & 0,8981 & 0,8066 & 0,8981 & 0,8499 \\
title      & WELFake & lstm + word2vec   & 0,8874 & 0,8162 & 0,8874 & 0,8467 \\
title      & ISOT    & cnn + word2vec    & 0,5434 & 0,8800 & 0,5434 & 0,6274 \\
short text & LIAR    & gru + word2vec    & 0,8977 & 0,8090 & 0,8977 & 0,8511 \\
short text & WELFake & lstm + word2vec   & 0,8286 & 0,8158 & 0,8286 & 0,8221 \\
short text & ISOT    & cnn + word2vec    & 0,6321 & 0,8800 & 0,6321 & 0,7050 \\
\bottomrule
\end{tabular}
\end{table}

W~ewaluacji na danych WWW sieci głębokie wykazują ciekawą cechę:
w~przeciwieństwie do klasycznych modeli najlepiej generalizują modele
wytrenowane na LIAR i~WELFake, mimo że in-domain plasowały się one
najniżej. Najlepszy wynik w~całej ewaluacji web osiąga BiLSTM~+~GloVe
trenowany na LIAR (\(F_1 = 0{,}8550\)) na pełnych tekstach
artykułów -- jest to wynik, który równa się najlepszym klasycznym
modelom i~wyprzedza transformery o~kilka punktów procentowych. Z~kolei
modele wytrenowane na ISOT istotnie tracą jakość na danych
zewnętrznych (\(F_1\) zaledwie \(0{,}63\)--\(0{,}71\)), co potwierdza
hipotezę, że korpus ISOT zawiera silne korelacje stylistyczne
(prefiks ,,WASHINGTON (Reuters)'', stałe sygnatury) wyciekające do
zbioru testowego.

\subsection{Podsumowanie}
\label{sec:wyniki_glebokie_podsum}

Sieci głębokie stanowią naturalny krok pośredni między klasycznym
uczeniem maszynowym a~modelami transformerowymi. Na korpusach ISOT
i~WELFake przekraczają próg \(F_1 = 0{,}97\), a~najlepsza architektura
(GRU~+~Word2Vec) jest niemal nieodróżnialna od klasyfikatora
XGBoost~+~TF-IDF. Na trudnym zbiorze LIAR ujawniają się natomiast
fundamentalne ograniczenia tej rodziny -- modele rekurencyjne
o~nieprzemyślanej konfiguracji mogą przewidywać wyłącznie klasę
większościową. Z~drugiej strony właśnie sieci głębokie najlepiej
generalizują na świeże dane WWW, co czyni je dobrym kandydatem do
zastosowań produkcyjnych. Stanowi to interesujący kontrast: model,
który nie uzyskał najwyższej miary $F_1$ na zbiorze testowym, może
okazać się najbardziej użyteczny w~rzeczywistej eksploatacji.

\section{Wyniki modeli opartych na transformerach}
\label{sec:wyniki_transformery}

W~rozdziale tym zaprezentowano wyniki czterech wstępnie wytrenowanych
modeli językowych z~rodziny Transformer: BERT (\textit{bert-base-uncased}),
DistilBERT (\textit{distilbert-base-uncased}), RoBERTa (\textit{roberta-base})
oraz FakeBERT -- wariant domenowy oparty na \textit{bert-base-uncased}
z~dodatkową adaptacją na korpusie artykułów newsowych. Modele zostały
dostrojone (fine-tuning) na każdym z~trzech zbiorów oddzielnie,
z~wykorzystaniem optymalizatora AdamW oraz schedulera \textit{linear
warmup}. Liczba kroków dostrajania, learning rate oraz rozmiar partii
zostały wybrane przez Optuna z~mechanizmem przedwczesnego zatrzymywania
(\textit{MedianPruner}).

\subsection{Wyniki na zbiorach testowych}
\label{sec:wyniki_bert_test}

\begin{table}[ht]
\centering
\caption{Wyniki modeli transformerowych na zbiorze testowym ISOT.}
\label{tab:bert_isot}
\begin{tabular}{llrrrrrr}
\toprule
Model & Tokenizer & Acc. & Prec. & Recall & $F_1$ & $t_{tr}$ [s] & $t_{inf}$ [s] \\
\midrule
fakebert   & bert-base-uncased       & 1,0000 & 1,0000 & 1,0000 & 1,0000 & 302,70 & 3,54 \\
roberta    & roberta-base            & 0,9995 & 0,9995 & 0,9995 & 0,9995 & 179,61 & 1,81 \\
bert       & bert-base-uncased       & 0,9990 & 0,9990 & 0,9990 & 0,9990 & 264,51 & 3,03 \\
distilbert & distilbert-base-uncased & 0,9982 & 0,9982 & 0,9982 & 0,9982 & 106,44 & 1,42 \\
\bottomrule
\end{tabular}
\end{table}

Na zbiorze ISOT wszystkie modele transformerowe osiągnęły wynik bardzo
bliski idealnemu. Spektakularnie wypadł FakeBERT, który dla 3911
próbek testowych nie pomylił się ani razu (macierz pomyłek
$\begin{pmatrix}1791 & 0\\0 & 2120\end{pmatrix}$). Również pozostałe
modele przekroczyły barierę \(F_1 = 0{,}998\). Warto jednak zauważyć,
że na zbiorze ISOT różnica między rodziną klasyczną a~transformerową
sięga jedynie \(0{,}2\)--\(0{,}3\) punktu procentowego $F_1$, a~koszt
obliczeniowy transformerów jest około 10-krotnie wyższy. Powstaje
zatem pytanie, czy nakład obliczeniowy jest uzasadniony.

\begin{table}[ht]
\centering
\caption{Wyniki modeli transformerowych na zbiorze testowym WELFake.}
\label{tab:bert_welfake}
\begin{tabular}{llrrrrrr}
\toprule
Model & Tokenizer & Acc. & Prec. & Recall & $F_1$ & $t_{tr}$ [s] & $t_{inf}$ [s] \\
\midrule
roberta    & roberta-base            & 0,9978 & 0,9978 & 0,9978 & 0,9978 & 702,81 & 5,79  \\
bert       & bert-base-uncased       & 0,9951 & 0,9951 & 0,9951 & 0,9951 & 563,95 & 5,54  \\
fakebert   & bert-base-uncased       & 0,9950 & 0,9950 & 0,9950 & 0,9950 & 856,57 & 10,41 \\
distilbert & distilbert-base-uncased & 0,9943 & 0,9944 & 0,9943 & 0,9943 & 286,44 & 3,64  \\
\bottomrule
\end{tabular}
\end{table}

Bardziej wymagający korpus WELFake pokazuje, że rodzina transformerów
oferuje wymierną przewagę nad klasycznymi modelami i~sieciami
głębokimi. RoBERTa osiągnęła tu najlepszy wynik (\(F_1 = 0{,}9978\)),
co stanowi zysk około \(2{,}6\) punktu procentowego względem
najlepszego klasycznego modelu (XGBoost~+~BoW) oraz \(1{,}5\) punktu
względem GRU~+~Word2Vec. Wszystkie cztery modele transformerowe
mieszczą się w~bardzo wąskim paśmie wyników (\(0{,}9943\)--\(0{,}9978\)),
co dowodzi spójności podejścia opartego na uwadze.

\begin{table}[ht]
\centering
\caption{Wyniki modeli transformerowych na zbiorze testowym LIAR.}
\label{tab:bert_liar}
\begin{tabular}{llrrrrrr}
\toprule
Model & Tokenizer & Acc. & Prec. & Recall & $F_1$ & $t_{tr}$ [s] & $t_{inf}$ [s] \\
\midrule
roberta    & roberta-base            & 0,6578 & 0,6560 & 0,6578 & 0,6564 & 50,48  & 0,400 \\
bert       & bert-base-uncased       & 0,6531 & 0,6537 & 0,6531 & 0,6534 & 31,98  & 0,685 \\
distilbert & distilbert-base-uncased & 0,6421 & 0,6434 & 0,6421 & 0,6426 & 27,95  & 0,263 \\
fakebert   & bert-base-uncased       & 0,6406 & 0,6380 & 0,6406 & 0,6383 & 285,38 & 1,349 \\
\bottomrule
\end{tabular}
\end{table}

Na zbiorze LIAR transformery uzyskały najlepsze wyniki w~całym
badaniu (\(F_1 = 0{,}6564\) dla RoBERTa) -- choć przewaga nad
najlepszym klasycznym modelem (MNB~+~BoW, \(F_1 = 0{,}6326\)) jest
umiarkowana (około \(2{,}4\) punktu procentowego), to jest istotnie
większa niż w~przypadku sieci głębokich, które na tym zbiorze
nie potrafią konsekwentnie pokonać klasyfikatora większościowego.
Z~drugiej strony FakeBERT, mimo długiego czasu uczenia, okazuje się
najsłabszym z~transformerów -- dodatkowy etap dostrajania
na artykułach nie pomaga w~klasyfikacji krótkich wypowiedzi
politycznych, a~wręcz utrudnia adaptację do tej domeny.

\subsection{Analiza czasów uczenia i~wnioskowania}
\label{sec:wyniki_bert_czas}

\begin{table}[ht]
\centering
\caption{Średnie czasy uczenia i~wnioskowania transformerów
w~zależności od zbioru danych (w~sekundach).}
\label{tab:bert_czasy}
\begin{tabular}{lrrrr}
\toprule
Zbiór & $\overline{t_{tr}}$ & $\max{t_{tr}}$ & $\overline{t_{inf}}$ & $\max{t_{inf}}$ \\
\midrule
ISOT    & 213,32 & 302,70 & 2,453 & 3,544  \\
LIAR    &  98,95 & 285,38 & 0,674 & 1,349  \\
WELFake & 602,44 & 856,57 & 6,346 & 10,411 \\
\bottomrule
\end{tabular}
\end{table}

Modele transformerowe są zdecydowanie najbardziej kosztowne pod
względem zasobów obliczeniowych. Średni czas uczenia na zbiorze
WELFake wynosi \(602{,}4\)~s, czyli około 18 razy więcej niż
w~przypadku rodziny klasycznej i~6 razy więcej niż sieci głębokich.
Pojedyncze wnioskowanie zajmuje 1--10~ms na próbkę, co oznacza,
że w~scenariuszu wsadowym (batch~=~32) możliwe jest osiągnięcie
przepustowości około 200~próbek/s na pojedynczej karcie GPU klasy
T4. Najszybszym transformerem pozostaje DistilBERT, oferując około
\(40\%\) krótszy czas inferencji w~stosunku do BERT przy zachowaniu
zbliżonej skuteczności.

\subsection{Wyniki na danych pozyskanych ze stron WWW}
\label{sec:wyniki_bert_web}

\begin{table}[ht]
\centering
\caption{Najlepsze modele transformerowe w~ewaluacji na danych WWW
(według miary $F_1$).}
\label{tab:bert_web}
\begin{tabular}{lllrrrr}
\toprule
Wariant & Zbiór & Model & Acc. & Prec. & Recall & $F_1$ \\
\midrule
text       & ISOT    & roberta    & 0,7311 & 0,8288 & 0,7311 & 0,7724 \\
text       & WELFake & fakebert   & 0,6629 & 0,7952 & 0,6629 & 0,7205 \\
text       & LIAR    & fakebert   & 0,5641 & 0,8247 & 0,5641 & 0,6498 \\
title      & ISOT    & roberta    & 0,7465 & 0,8630 & 0,7465 & 0,7899 \\
title      & WELFake & fakebert   & 0,7547 & 0,7960 & 0,7547 & 0,7747 \\
title      & LIAR    & fakebert   & 0,2711 & 0,7905 & 0,2711 & 0,3273 \\
short text & ISOT    & roberta    & 0,8418 & 0,8741 & 0,8418 & 0,8556 \\
short text & ISOT    & fakebert   & 0,7571 & 0,8909 & 0,7571 & 0,8017 \\
short text & LIAR    & fakebert   & 0,5191 & 0,8622 & 0,5191 & 0,6078 \\
\bottomrule
\end{tabular}
\end{table}

W~ewaluacji na danych pozyskanych ze stron WWW transformery
prezentują niejednoznaczny obraz. Z~jednej strony osiągają najwyższą
średnią miarę $F_1$ wśród wszystkich rodzin dla pełnych tekstów
(\(\overline{F_1} = 0{,}651\) wobec \(0{,}580\) dla klasycznych
i~\(0{,}576\) dla sieci głębokich), co świadczy o~ich potencjale do
generalizacji. Z~drugiej strony w~scenariuszu samych nagłówków
i~krótkich fragmentów ulegają sieciom głębokim wytrenowanym na LIAR.
Najlepszą pojedynczą konfiguracją jest RoBERTa wytrenowana na ISOT,
która na krótkich fragmentach uzyskuje \(F_1 = 0{,}8556\) -- jest to
najwyższa miara $F_1$ ewaluacji web-scrapingu w~podziale na typ
testu \textit{short text}.

\subsection{Podsumowanie}
\label{sec:wyniki_bert_podsum}

Modele oparte na architekturze Transformer osiągają najwyższe wyniki
w~zadaniu detekcji fałszywych informacji we wszystkich trzech badanych
zbiorach. Na ISOT idealną skuteczność uzyskuje FakeBERT (\(F_1 = 1{,}0\)),
na WELFake -- RoBERTa (\(F_1 = 0{,}9978\)), a~na trudnym zbiorze LIAR
ponownie RoBERTa (\(F_1 = 0{,}6564\)), choć w~tym przypadku różnica
względem prostych modeli klasycznych jest niewielka. Ich główną
słabością jest wysoki koszt obliczeniowy -- średnio kilkukrotnie
wyższy zarówno na etapie uczenia, jak i~wnioskowania, niż w~przypadku
pozostałych rodzin. Dodatkowo na danych spoza zbioru treningowego
przewaga transformerów nad sieciami głębokimi jest niewielka lub
wręcz odwraca się na korzyść tych drugich. Daje to ciekawą wskazówkę
dla projektowania systemów produkcyjnych: w~scenariuszach, w~których
dystrybucja danych testowych różni się od treningowej, prostsze
modele potrafią dorównać, a~niekiedy nawet przewyższyć kosztowne
transformery.

% =====================================================================
% Koniec rozdziałów z~wynikami eksperymentów.
% =====================================================================
